\section{Proof of the logarithmic-tube estimates}
\label{app:growing-tube}

We prove the growing clauses of \cref{thm:history-graph}.  Throughout this
appendix, a target \(\del\) is fixed,
\(S=\widehat S_\del=S_\del+B_{\rm hit}\), and all cutoffs are frozen.
Taylor differentiation is with respect to an independent amplitude
\(\rho\in[-\del,\del]\); in particular, neither \(\widehat S_\del\) nor a
cutoff profile is differentiated.

\subsection{Canard flow and nested tubes}

In the coordinates \cref{eq:canard-coordinates}, the singular flow is
\cref{eq:canard-coordinate-flow}.  If a backward segment has fixed length
\(0\le t\le T\), then its normal linear factor on \(d=0\) is
\[
 \exp\left(-\chi t+\frac{t^2}{2}\right).
\]
Consequently, for every fixed derivative order \(k\), there are
\(C_{k,T},c_{k,T},m_k\) such that
\begin{equation}
\label{eq:fixed-flow-jet}
 \bigl|D^k\phi_0^{-t}(\Gamma(\chi,d))\bigr|
 \le C_{k,T}\langle\chi\rangle^{m_k}\e^{c_{k,T}|\chi|}
\end{equation}
while the segment remains in a fixed-width normal tube.  To see this,
observe first that \(|\chi(t)|\le|\chi(0)|+C_T\); Gronwall proves the
first-variation bound, and the higher variational equations have the same
homogeneous part with polynomial sources in lower variations.

Let \(J\) be the largest grade in the finite rectangular jet set.  Order
the jet blocks as in \cref{app:mixed-jets} and denote their number by
\(N_{\rm bl}\).  Take
\[
 D=2N_{\rm bl}+4,\qquad
 1>\kappa_0>\kappa_1>\cdots>\kappa_D>0,
\]
put \(T_*=\theta_1+1\), \(B_j=jT_*\), and use the outer buffer
\[
 B_*=B_D+T_*+2.
\]
Define
\begin{equation}
\label{eq:nested-tubes}
 \cU_j(\del)=\left\{\Gamma(\chi,d):
 |\chi|\le S+B_j,\ |d|\le\del^{\kappa_j}\right\}.
\end{equation}
The cutoff equals one on the fixed-width flow hull of
\(|\chi|\le S+B_*\) and is supported in a twice larger longitudinal
interval.  Its finite derivatives and those of the prepared polynomial
data are bounded by polynomials \(P_g(S)\), \(g\le J+1\).  The cutoff
field is complete.

\subsection{Global contraction and the stable tail}

For every fixed jet grade, \cref{eq:fixed-flow-jet} and the cutoff bounds
give a global composition bound
\begin{equation}
\label{eq:global-Bg}
 B_g(S)=P_g(S)\e^{C_gS}.
\end{equation}
Since
\[
 \del P_g(S)\e^{C_gS}\longrightarrow0,
\]
the order-zero transform and every highest jet block have contraction
factor
\begin{equation}
\label{eq:growing-contraction}
 q_\del\le\del P_J(S)\e^{C_JS}=o(1).
\end{equation}
The common-fiber argument in \cref{app:mixed-jets}, applied with these
target-dependent radii, therefore gives a unique global fixed point for
\(|\rho|\le\del\), its finite \(C_\rho^3C_\nu^1C_\eta^2\) jets, and the
global bounds \cref{eq:global-Bg}.

For the stable convolution, set
\[
 R_{{\rm tail},\del}=L\log(1/\del)
\]
and split the full integral in \cref{eq:graph-transform} at
\(R_{{\rm tail},\del}\).  The undifferentiated tail is bounded by
\[
 \frac{M_A}{\beta}|\rho|\|G\|_\infty\del^{\beta L}.
\]
After any of the finitely many differentiations, flow jets add only a fixed
power of \(r\), and half of the semigroup decay can be retained:
write \(\mathscr D\) for any derivative in the required rectangular
\((u,\rho,\nu,\eta)\)-jet family.  Then
\begin{equation}
\label{eq:stable-tail}
 \bigl\|\mathscr D\cT_H^{>R_{{\rm tail},\del}}\bigr\|_\infty
 \le P_J(S)\e^{C_JS}(1+R_{{\rm tail},\del})^{M_J}
      \e^{-\beta R_{{\rm tail},\del}/2}.
\end{equation}
Choose \(L\) after \(J\) so that the tail is \(O(\del^6)\), three orders
beyond the leading graph coefficient retained in the main theorem.  Since all
prefactors in \cref{eq:stable-tail} are \(\del^{-o(1)}\), the tail is
\(O(\del^6)\) for the entire rectangular family.  This splitting is only
an estimate: the fixed point still uses the exact integral over
\([0,\infty)\), so invariance is not truncated.

For \(0\le r\le R_{{\rm tail},\del}\), the additional
stable-convolution backtrack has
length
\[
 |\rho|r\le\del R_{{\rm tail},\del}=o(1).
\]
After decreasing \(\del_0\), \cref{eq:fixed-flow-jet} and the gaps between
the \(\kappa_j\) imply
\begin{equation}
\label{eq:nested-inclusion}
 \del^{\kappa_j}P_J(S)\e^{C_JS}
 +\del P_J(S)\e^{C_JS}
 \le\frac12\del^{\kappa_{j+1}}.
\end{equation}
Thus a current or delayed evaluation starting in \(\cU_j\) lies in
\(\cU_{j+1}\); at the terminal level it remains in the uncut buffer.

\subsection{Local mixed-jet fibers}

Let \(\cT_{\rho,S}\) denote the graph transform with every cutoff frozen at
the target \(S\) and with the perturbation amplitude replaced by \(\rho\).
Let \(q_{0,\del}^{\rm pr}\) be the complete singular preparation fixed by
\(\cP\) at this target.

For a jet tensor \(V\), define
\[
 \|V\|_{j;c,m}
 =\sup_{u=\Gamma(\chi,d)\in\cU_j}
 \frac{|V(u)|}{\langle\chi\rangle^m\e^{c|\chi|}}.
\]
Let \(Z_{n+1}=\cT_{\rho,S}Z_n\),
\(Z_0=(q_{0,\del}^{\rm pr},0)\), and let
\(\mathbf J_\ell\) denote the \(\ell\)-th ordered mixed-jet block.  The
Faà di Bruno expansion of one graph-transform step has the exact
highest-block structure
\begin{equation}
\label{eq:local-highest-block}
 \mathbf J_\ell Z_{n+1}
 =b_\ell(\mathbf J_{<\ell}Z_n)
  +\rho\mathscr L_{\ell,n}[\mathbf J_\ell Z_n]
  +E_{\ell,n}^{>R_{{\rm tail},\del}}.
\end{equation}
Here \(E_{\ell,n}^{>R_{{\rm tail},\del}}\) is the differentiated stable-tail
term from \(r>R_{{\rm tail},\del}\).  The lower-block maps below are fixed
finite polynomials,
denoted \(\Pi_{\ell,h}\), with coefficients controlled by the prepared data.
For the first pure spatial block, the current highest-block map is nonlinear
rather than affine, but its difference estimate has the same factor
\(|\rho|\).  If a \(\rho\)-derivative removes the displayed factor, the
remaining graph jet has one lower \(\rho\)-grade and belongs to
\(b_\ell\).  Every other factor in \(b_\ell\) is a lower graph jet or a
flow jet evaluated one nested level farther out.

Set
\[
 A_{\ell,j}^{(n)}
 =\|\mathbf J_\ell Z_n\|_{j;c_\ell,m_\ell}.
\]
After enlarging the finite weights, \cref{eq:local-highest-block} gives
\begin{equation}
\label{eq:local-recurrence}
 A_{\ell,j}^{(n+1)}
 \le C_\ell+
 C_\ell\sum_{h<\ell}\Pi_{\ell,h}
      \bigl(A_{h,j+1}^{(n)}\bigr)
 +|\rho|P_\ell(S)\e^{C_\ell S}B_\ell(S)
 +C_\ell\del^6 .
\end{equation}
Induction over \(\ell\) proves
\begin{equation}
\label{eq:local-induction}
 \sup_nA_{\ell,j}^{(n)}\le C_\ell,
 \qquad 0\le j\le D-\ell.
\end{equation}
Indeed, the assertion is immediate for the zeroth block.  If it holds for
\(h<\ell\), then \(j+1\le D-h\) bounds every lower-block term in
\cref{eq:local-recurrence}; the sole global current-block term is
\(o(1)\) by \cref{eq:growing-contraction}.  Since
\(D>N_{\rm bl}\), the core \(j=0\) is covered for every required block.
Global convergence of the jets implies convergence in each local seminorm,
so \cref{eq:local-induction} passes to the fixed point.

The induction closes because an \(S\)-dependent lower jet is always
controlled in the next weighted local fiber, whereas the only use of a
global current jet retains \(\rho\).  Infinite nesting in the exact fixed
point is summed by the global contraction and does not create infinite atom
depth in a finite Taylor coefficient.

\subsection{Atom depth and the Taylor remainder}

For a set \(\cU\), define the depth-\(d\) continuous flow hull
\[
 \mathfrak H^{[d]}(\cU)
 =\{\phi_0^{-t}u:u\in\cU,\ 0\le t\le d\theta_1\}.
\]
At \(\rho=0\), one differentiation of the transform gives
\[
 Q_1=F(\cE_0;0,\nu,\eta),\qquad
 H_1=-A^{-1}G(\cE_0;0,\nu,\eta),
\]
so \((Q_1,H_1)\) has depth one.  A second differentiation introduces at
most one additional delayed evaluation, one first variation
\[
 \int_0^{\theta_j}
 D\phi_0^{-(\theta_j-a)}
 Q_1(\phi_0^{-a}u)\,\dd a,
\]
or a moment of the stable kernel.  Hence \((Q_2,H_2)\) depends only on
\(\mathfrak H^{[2]}(\cU_0(\del))\).  The preparation equals the
physical field on that hull; these coefficients are therefore independent
of the admissible outer cutoff.

Apply \cref{eq:local-induction} to
\[
 D_u^a\partial_\nu^i\partial_\eta^j
 \partial_\rho^3(Q_{\rho,S},H_{\rho,S}),
 \qquad a\le3,\quad i\le1,\quad j\le2.
\]
Taylor's integral formula in the fixed common fiber yields
\begin{align}
 Q_{\del,S}&=q_0+\del Q_1+\del^2Q_2+\del^3R_{3,Q},\\
 H_{\del,S}&=\del H_1+\del^2H_2+\del^3R_{3,H},
 \label{eq:appendix-growing-expansion}
\end{align}
where
\begin{equation}
\label{eq:appendix-growing-remainder}
 \max_{a\le3,\ i\le1,\ j\le2}
 \bigl|D_u^a\partial_\nu^i\partial_\eta^j
      (R_{3,Q},R_{3,H})(\Gamma(\chi,d))\bigr|
 \le C\langle\chi\rangle^m\e^{c|\chi|}
\end{equation}
on \cref{eq:core-tube}.  This is the enlarged target used by the moving-hit
construction from the outset, not a second graph fixed point.  Because
\(B_{\rm hit}\) is fixed in \(\cP\), every factor
\(P(\widehat S_\del)\e^{C\widehat S_\del}\) is bounded by a constant times
the corresponding factor with \(S_\del\), and
\(\del P(\widehat S_\del)\e^{C\widehat S_\del}=o(1)\).  Hence the delayed
evaluations from this enlarged current tube remain in the farther nested
uncut hull with the same estimates.

The Gaussian weight in
\cref{eq:first-integral-gradient} makes
\cref{eq:appendix-growing-remainder} integrable with no
\(\del^{-o(1)}\) loss.

Finally, extend the boundary state jets through degree three in the
unscaled normal distance and multiply by a fixed-width cutoff.  No factor
\(\del^{-\kappa_0}\) is generated.  The prepared planar field has the
same rectangular bounds and equals the actual graph field on the shrinking
tube.  With target \(\widehat S_\del\), the bootstrap
\cref{eq:normal-bootstrap} puts the retained current points inside the
shrinking normal component of the enlarged tube and their depth-two
backtracks inside the nested uncut hull.  Therefore the full
Lyapunov--Perron fixed point, rather than a truncated approximation, supplies
the exact embedding \cref{eq:history-embedding} and semiconjugacy
\cref{eq:exact-semiconjugacy} on every retained trace.  This completes the
proof of the growing clauses of \cref{thm:history-graph}.
